\section{Introducción}

El contactor electromagnético representa uno de los dispositivos fundamentales en el control y automatización de sistemas eléctricos industriales, actuando como el elemento de maniobra encargado de interrumpir o establecer el flujo de potencia hacia cargas de gran demanda. Esta práctica tiene como finalidad el estudio profundo del comportamiento del circuito magnético del contactor, analizando las variables eléctricas críticas, como la corriente de llamada y de mantenimiento, así como su respuesta bajo condiciones de operación normal. A través del desarmado del equipo y la ejecución de mediciones experimentales, se busca contrastar el funcionamiento práctico del dispositivo con los lineamientos técnicos establecidos por el fabricante y las normas internacionales vigentes, permitiendo así comprender la relación funcional entre su estructura mecánica y sus parámetros eléctricos operativos.